\documentclass[11pt]{article}
\usepackage[T1]{fontenc}
\usepackage[margin=1in]{geometry}
\usepackage{amsmath,amssymb,amsthm}
\usepackage{url}

\newtheorem{theorem}{Theorem}
\newtheorem{lemma}{Lemma}
\newtheorem{proposition}{Proposition}
\newtheorem{corollary}{Corollary}
\theoremstyle{definition}
\newtheorem{definition}{Definition}
\newtheorem{hypothesis}{Hypothesis}
\theoremstyle{remark}
\newtheorem{remark}{Remark}

\newcommand{\PHP}{\mathrm{PHP}}
\newcommand{\Fub}{\mathrm{Fub}}
\newcommand{\pig}{\mathrm{pig}}
\newcommand{\mex}{\operatorname{mex}}
\newcommand{\UP}{\mathrm{UP}}
\newcommand{\todo}[1]{\textbf{[TODO: #1]}}
\newcommand{\N}{\mathbb{N}}

\title{Exact conflict counts of unit-propagating DPLL on the pigeonhole principle:\\
closed forms for pigeon-major and hole-major static orders}
\author{Sai Arun Kumar Katherashala}
\date{October 2026}

\begin{document}
\maketitle

\begin{abstract}
We study the number of conflict leaves of the DPLL procedure with unit propagation on the
pigeonhole formula $\PHP(n{+}1,n)$ (pigeon clauses and hole-pair clauses only) when the
branching variable is chosen by a fixed static order of the $n(n{+}1)$ variables. For the
\emph{pigeon-major} order we prove that the number of conflict leaves is $n\cdot\Fub(n{-}1)$, where
$\Fub$ denotes the ordered Bell (Fubini) numbers (OEIS A052882). For the \emph{hole-major} order
we prove that it is $(n+2)^{n-2}+(n+3)^{n-2}$ for $n\ge 2$. We derive the exponential generating
functions and the resulting asymptotics. We also report exact values for the anti-diagonal static
order ($n\le 16$; the value for $n=16$ is from a single run, see Section~\ref{sec:limits}) and for a dynamic Jeroslow--Wang rule with 4-way branching (``$\mathcal D_2$-JW'', called Res(2)-JW in earlier drafts and not Krajicek's Res(2),
$n\le 14$), for which we know no closed form, and we prove a conditional statement about the
$\mathcal D_2$-JW generating function (conditional on an unproved growth hypothesis). These counts are
properties of a branching heuristic; they are \emph{not} resolution proof sizes, and nothing here
is a lower bound for any proof system. We also list the claims of earlier drafts of this work that
are not established or are false, with counterexamples.
\end{abstract}

%=====================================================================
\section{Introduction}\label{sec:intro}

\paragraph{Setting.}
Let $\PHP(n{+}1,n)$ be the CNF formula over the variables $x_{i,j}$
($i\in\{0,\dots,n\}$ a pigeon, $j\in\{1,\dots,n\}$ a hole) consisting of the $n{+}1$
\emph{pigeon clauses} $P_i=\bigvee_{j=1}^n x_{i,j}$ and the $n\binom{n+1}{2}$ \emph{hole clauses}
$(\neg x_{i,j}\vee\neg x_{k,j})$, $0\le i<k\le n$, $1\le j\le n$. There are no functional
clauses (a pigeon may occupy several holes) and no onto clauses; this is the formulation used by
Haken~\cite{Haken85}. The formula is unsatisfiable.

\paragraph{The procedure.}
\emph{DPLL with unit propagation} (we write $\UP$) is the following search procedure
\cite{DP60,DLL62}. At a node with a partial assignment $\alpha$, run unit propagation: while some
clause has no true literal and exactly one unassigned literal, make that literal true; stop with a
\emph{conflict} if some clause has all literals false. If there is a conflict the node is a
\emph{conflict leaf}. Otherwise let $x$ be the first unassigned variable in a fixed \emph{static
order} and create two children, $\alpha\cup\{x{=}1\}$ and $\alpha\cup\{x{=}0\}$. Only unit clauses
are used (no pure literals, no learning). Because $\PHP$ is unsatisfiable every maximal branch ends
in a conflict, so the search tree is a full binary tree whose number of nodes is twice its number of
conflict leaves minus one. The number of conflict leaves does not depend on which child is explored
first (the tree is the set of all children of every node; the branching variable of a node depends only on the set of assigned variables for a static order, and on the assignment itself for JW, never on the order in which children are visited). Let $C_\prec(n)$ denote the number of conflict leaves for the static order $\prec$.

\paragraph{Orders.} In terms of the variable $x_{i,j}$:
\begin{itemize}
\item \emph{pigeon-major}: $x_{i,j}\prec x_{i',j'}$ iff $(i,j)<(i',j')$ lexicographically
  (row by row);
\item \emph{hole-major}: iff $(j,i)<(j',i')$ lexicographically (column by column, each column from
  pigeon $0$ down to pigeon $n$);
\item \emph{anti-diagonal}: iff $(i+j,\,i)<(i'+j',\,i')$ lexicographically (anti-diagonals
  $i+j=\text{const}$, increasing pigeon inside each anti-diagonal).
\end{itemize}

\paragraph{What these numbers are, and what they are not.}
$C_\prec(n)$ is a property of a branching heuristic on a fixed formula. It is \emph{not} the size of
a resolution refutation. A DPLL tree with unit propagation turns into a tree-like resolution
refutation only after each conflict leaf is expanded into a unit-resolution derivation, which adds
leaves. For instance, $\PHP(2,1)$ has $C(1)=1$ in every order (the root already propagates to a
conflict), but every tree-like resolution refutation of it has $3$ leaves (the axioms $x_{0,1}$,
$x_{1,1}$, $\neg x_{0,1}\vee\neg x_{1,1}$). The minimum number of leaves of a decision tree whose
every leaf falsifies an axiom (equivalently, the minimum size in leaves of a tree-like resolution
refutation) is $3,11,43,189$ for $n=1,2,3,4$ (Section~\ref{sec:comp}); the pigeon-major counts
$1,2,9,52$ are smaller. General resolution is governed by Haken's theorem \cite{Haken85}:
every resolution refutation of $\PHP(n{+}1,n)$ has size $2^{\Omega(n)}$ (see also
\cite{CR79} for the framework of proof systems). \emph{We make no claim about lower bounds for
any proof system, and none of our results should be read as one.}

\paragraph{Results.}
\begin{enumerate}
\item[(A)] (Theorem~\ref{thm:pig}) Pigeon-major: $C(n)=n\cdot\Fub(n-1)$ for all $n\ge1$.
\item[(B)] (Theorem~\ref{thm:hole}) Hole-major: $C(1)=1$ and $C(n)=(n+2)^{n-2}+(n+3)^{n-2}$ for
  all $n\ge2$.
\item[(C)] Exact values for the anti-diagonal order ($n\le16$; $n=16$ from a single run) and for the $\mathcal D_2$-JW rule
  ($n\le14$), for which no closed form is known to us (Section~\ref{sec:comp}).
\item[(D)] Generating functions and asymptotics for (A) and (B) (Section~\ref{sec:asym}).
\item[(E)] A conditional theorem on the generating function of the $\mathcal D_2$-JW counts
  (Theorem~\ref{thm:C}, Section~\ref{sec:cond}), resting on an explicitly stated, \emph{unproved}
  Hypothesis~\ref{hyp:H}.
\item[(G)] The congruence $r(n)\equiv1\pmod 3$ for the $\mathcal D_2$-JW counts (Proposition~\ref{prop:mod3}).
\item[(F)] A list of claims in earlier drafts of this work that are false or not established
  (Section~\ref{sec:errata}).
\end{enumerate}
To our knowledge the closed form (B) and the generating-function form of the hole-major counts
are new; (A) gives a combinatorial proof for a sequence that coincides with A052882 \cite{OEIS}.
We have made no systematic literature search beyond Section~\ref{sec:related}.

\paragraph{Reproducibility.} Every numerical statement in this paper is re-verified by the script
\texttt{ms\_check.py} (Section~\ref{sec:repro}).

%=====================================================================
\section{Preliminaries}\label{sec:prelim}

A partial assignment $\beta$ is \emph{closed} (for a CNF $F$) if no clause of $F$ is falsified and no
clause is unit (no true literal and exactly one unassigned literal). A sequence of literals
$\ell_1,\dots,\ell_k$ is \emph{forced from} $\alpha$ if for each $s$ there is a clause $C_s\in F$
all of whose literals except $\ell_s$ are false under $\alpha_{s-1}:=\alpha\cup\{\ell_1,\dots,\ell_{s-1}\}$
and $\ell_s$ is unassigned under $\alpha_{s-1}$. A unit-propagation run from $\alpha$ is a maximal
forced sequence; it ends either with a falsified clause (conflict) or at a closed assignment.

\begin{lemma}[UP is confluent]\label{lem:conf}
Let $F$ be a CNF and $\alpha$ a partial assignment.
\begin{enumerate}
\item If $\beta\supseteq\alpha$ is closed, then every forced sequence from $\alpha$ is contained in
  $\beta$ and the assignment $\alpha\cup\{\ell_1,\dots,\ell_k\}$ falsifies no clause.
\item Consequently, unit propagation from $\alpha$ ends in a conflict (for every run) iff $\alpha$
  has no closed extension; otherwise all runs end in the same closed assignment $\UP(\alpha)$,
  the least closed extension of $\alpha$.
\item If $\gamma=\alpha\cup\{\ell_1,\dots,\ell_k\}$ for a forced sequence and $\gamma$ is closed, then
  $\UP(\alpha)=\gamma$.
\end{enumerate}
\end{lemma}
\begin{proof}
(1) Induction on $s$: if $\ell_1,\dots,\ell_{s-1}\in\beta$, the clause $C_s$ has all literals but
$\ell_s$ false under $\beta$; as $\beta$ is closed, $C_s$ is neither falsified nor unit, so $\ell_s$
is true in $\beta$. A clause falsified by $\alpha\cup\{\ell_1,\dots,\ell_k\}\subseteq\beta$ would be
falsified by $\beta$.
(2) If $\beta$ is a closed extension, by (1) no run produces a conflict, so every run ends at a closed
assignment $\beta^*\subseteq\beta$. Two runs ending at $\beta^*_1,\beta^*_2$ satisfy
$\beta^*_1\subseteq\beta^*_2$ (apply this to the closed extension $\beta^*_2$) and vice versa, so they
coincide, and each is contained in every closed extension. If there is no closed extension, a run
cannot end without conflict (it would end at a closed extension).
(3) $\gamma$ is a closed extension, so $\UP(\alpha)\subseteq\gamma$ by (2); and $\gamma\subseteq\UP(\alpha)$
by (1) applied to $\beta=\UP(\alpha)$. \qedhere
\end{proof}

In particular, whether a node is a conflict leaf is independent of the order of propagation, the
children of a non-leaf node are $\UP(\UP(\alpha)\cup\{x{=}v\})$, and in the proofs below we may
propagate in any convenient order. Two consequences used repeatedly:
\begin{itemize}
\item[(U-i)] if we exhibit a closed assignment $\gamma$ reached from $\alpha$ by forced steps, then
  $\gamma$ \emph{is} the result of unit propagation;
\item[(U-ii)] if no closed extension of $\alpha$ exists (e.g.\ every extension without a falsified
  clause leaves a unit clause or falsifies a hole clause), then unit propagation from $\alpha$ ends in
  a conflict.
\end{itemize}

\paragraph{Fubini numbers.} $\Fub(0)=1$ and $\Fub(m)=\sum_{k=1}^m\binom mk\Fub(m-k)$, the ordered Bell
numbers $1,1,3,13,75,541,\dots$ (OEIS A000670) with exponential generating function $1/(2-e^x)$.
The sequence $\mathrm{A052882}(n)=n\cdot\Fub(n-1)$ is OEIS A052882, with e.g.f.\ $x/(2-e^x)$.

%=====================================================================
\section{Pigeon-major order}\label{sec:pig}

\begin{theorem}[pigeon-major]\label{thm:pig}
For the pigeon-major order, $C(n)=n\cdot\Fub(n-1)$ for all $n\ge1$.
\end{theorem}

\paragraph{Case $n=1$.} The pigeon clauses are the units $x_{0,1}$ and $x_{1,1}$; propagation makes both
true and falsifies the hole clause. The root is a conflict leaf: $C(1)=1=1\cdot\Fub(0)$. From now on
$n\ge2$. The variables are ordered $x_{0,1},\dots,x_{0,n},x_{1,1},\dots,x_{n,n}$.

\subsection{States}
\emph{Row-start data} is $(i;T_0,\dots,T_{i-1})$ where the $T_p\subseteq[n]$ ($p<i$) are nonempty and
pairwise disjoint. Put $O_{i-1}=\bigcup_{p<i}T_p$ and let $c_1<\dots<c_f$ be the elements of
$[n]\setminus O_{i-1}$ (the \emph{free columns}), $f=n-|O_{i-1}|$. We only use data with $f\ge2$.
Since $|O_{i-1}|\ge i$ we then have $i\le n-f\le n-2$, so the pigeons $i,i+1,i+2$ exist.

For $0\le t\le f$ and $T\subseteq\{c_1,\dots,c_t\}$ put $O=O_{i-1}\cup T$ and define the partial
assignment $\sigma=\sigma(i;T_0,\dots,T_{i-1};t,T)$:
\begin{itemize}
\item[(S1)] for $p<i$: $x_{p,h}$ is true iff $h\in T_p$ (all assigned);
\item[(S2)] row $i$: $x_{i,h}$ is false for $h\in O_{i-1}$; for $k\le t$, $x_{i,c_k}$ is true iff
  $c_k\in T$; the variables $x_{i,c_k}$, $k>t$, are unassigned;
\item[(S3)] for $q>i$: $x_{q,h}$ is false for $h\in O$, and unassigned otherwise.
\end{itemize}
$\sigma$ is \emph{admissible} if $T\ne\emptyset$, or $T=\emptyset$ and $t\le f-2$. Put
$F(\sigma)=n-|O|=f-|T|$ (the number of columns not covered by true cells of rows $\le i$).
The \emph{row-start state} of row $i$ is $\sigma(i;\dots;0,\emptyset)$.

\begin{lemma}\label{lem:pig1}
Let $\sigma$ be admissible, $f\ge2$.
\begin{enumerate}
\item[(a)] If $F\ge1$ no clause is falsified by $\sigma$; if $F\ge2$ no clause is unit, i.e.\ $\sigma$ is
  closed.
\item[(b)] If $F\le1$, unit propagation from $\sigma$ ends in a conflict.
\end{enumerate}
\end{lemma}
\begin{proof}
Pigeon clauses. For $p<i$, $P_p$ contains the true literal $x_{p,h}$, $h\in T_p\ne\emptyset$.
$P_i$: if $T\neq\emptyset$ it is satisfied; if $T=\emptyset$, its literals $x_{i,h}$ with $h\in O_{i-1}$
and $x_{i,c_k}$, $k\le t$, are false and the $f-t\ge2$ literals $x_{i,c_k}$, $k>t$, are unassigned, so it
is neither falsified nor unit. For $q>i$, $P_q$ has false literals exactly in the columns of $O$ and
the $F$ others unassigned: it is falsified iff $F=0$ and unit iff $F=1$.
Hole clauses. Fix a column $h$. If $h\in O_{i-1}$ or $h\in T$, exactly one variable of the column is
true (disjointness of the $T_p$, and $x_{q,h}$ is false for $q>i$ by (S3)), so every hole clause of the column
contains a false variable (hence a true literal) and is satisfied. If $h\notin O$, no variable of the column is true;
rows $p<i$ are false, row $i$ is false or unassigned, rows $q>i$ are unassigned, so every hole clause
$\neg x_{p,h}\vee\neg x_{q,h}$ has a false variable (satisfied) or two unassigned variables (not unit).
This proves (a).
(b) If $F=0$, every literal of $P_{i+1}$ is false: conflict. If $F=1$ let $h_0$ be the unique column
outside $O$. Since $f\ge 2$, pigeons $i+1$ and $i+2$ exist, and in $P_{i+1}$, $P_{i+2}$ all literals
except $x_{i+1,h_0}$, $x_{i+2,h_0}$ are false. A closed extension of $\sigma$ would have to make both
$x_{i+1,h_0}$ and $x_{i+2,h_0}$ true (otherwise one of these pigeon clauses is falsified or unit), which
falsifies a hole clause. So no closed extension exists and (U-ii) gives a conflict.
\end{proof}

\begin{lemma}[transitions]\label{lem:pig2}
Let $\sigma$ be admissible with $F\ge2$ and $t<f$, and let $c=c_{t+1}$. The first unassigned
variable of $\sigma$ in pigeon-major order is $x_{i,c}$ (rows $<i$ are complete, and in row $i$ the cells in
$O_{i-1}$ and $c_1,\dots,c_t$ are assigned). The two children of the node $\sigma$ are:
\begin{enumerate}
\item $x_{i,c}=1$: unit propagation additionally sets $x_{q,c}=0$ for $q>i$, giving
  $\sigma'=\sigma(i;\dots;t{+}1,T\cup\{c\})$, admissible, $F'=F-1$.
\item $x_{i,c}=0$, and not ($T=\emptyset$ and $t+1=f-1$): the assignment
  $\sigma'=\sigma(i;\dots;t{+}1,T)$, admissible, $F'=F$.
\item $x_{i,c}=0$, $T=\emptyset$, $t+1=f-1$: the only unassigned literal of $P_i$ is $x_{i,c_f}$, so
  propagation sets it true and then $x_{q,c_f}=0$ for $q>i$, giving
  $\sigma'=\sigma(i;\dots;f,\{c_f\})$, admissible, $F'=f-1$.
\end{enumerate}
In every case the child is a conflict leaf iff $F'\le1$; if $F'\ge2$ its propagated assignment is exactly
$\sigma'$. If $\sigma'$ has $t'=f$ and $F'\ge2$ then $\sigma'$ equals the row-start state
$\sigma(i{+}1;T_0,\dots,T_{i-1},T';0,\emptyset)$ with $f'=F'\ge2$ free columns.
\end{lemma}
\begin{proof}
In each case $\sigma'$ is obtained from $\sigma$ plus the decision by forced steps (the hole clauses
$(\neg x_{i,c}\vee\neg x_{q,c})$ in case 1, the unit $P_i$ and then the same hole clauses in case 3;
$x_{p,c}$ for $p<i$ is already false). Admissibility is checked directly: in case 2, $T=\emptyset$
forces $t+1\le f-2$ because $t\le f-2$ and $t+1\ne f-1$. If $F'\ge2$, $\sigma'$ is closed by
Lemma~\ref{lem:pig1}(a), hence it is the propagation result by (U-i); if $F'\le1$ there is a conflict by
Lemma~\ref{lem:pig1}(b) (the forced literals are consequences of the node, so the conflict is
reached from the node). The identification with the row-start state of row $i{+}1$ follows by
comparing (S1)--(S3): row $i{+}1$ is false on $O_i=O_{i-1}\cup T'$ and unassigned elsewhere, as are
all later rows. A node with $T=\emptyset$ and $t=f-1$ never occurs by admissibility; this is why case 3
has no second branch.
\end{proof}

At the root ($n\ge2$) no clause is unit or falsified (pigeon clauses have $n\ge2$ literals, hole
clauses are binary), so the root is the row-start state $\sigma(0;\,;0,\emptyset)$ with $f=n$. Hence
every non-leaf node of the tree is an admissible state with $F\ge2$, and by Lemma~\ref{lem:pig2}
its children are determined as stated.

\subsection{The recurrence}
\begin{lemma}\label{lem:pig3}
Let $\nu$ be a row-start node with $f\ge2$ free columns. The number of conflict leaves below $\nu$ is
$L(f)$, depending only on $f$, where $L(1)=1$ and
\[
L(f)=f+\sum_{j=1}^{f-2}\binom fj L(f-j)\qquad(f\ge2).
\]
\end{lemma}
\begin{proof}
Strong induction on $f$. Inside row $i$ the transitions of Lemma~\ref{lem:pig2} depend only on
$f$, $t$, the size of $T$ and whether $T=\emptyset$. Encode a node reached in row $i$ by the word
$v_1v_2\cdots\in\{\mathsf T,\mathsf F\}^*$ where $v_k$ is the value of $x_{i,c_k}$; in case~3 the
forced value $v_f=\mathsf T$ is appended after $v_1\cdots v_{f-1}=\mathsf F^{f-1}$. Then $|T|$ is the
number of letters $\mathsf T$ and $F=f-|T|$, so by Lemma~\ref{lem:pig2} a node is a conflict leaf
iff its word has at least $f-1$ letters $\mathsf T$; since $|T|$ grows by at most one per letter, this
happens exactly when the number of $\mathsf T$ first equals $f-1$. The letters $v_1,\dots,v_{f-1}$ are
always genuine binary decisions, and $v_f$ is a genuine decision unless $v_1\cdots v_{f-1}=\mathsf F^{f-1}$,
in which case it is forced to $\mathsf T$. The word $\mathsf F^f$ never occurs.

\emph{Leaves in row $i$.} The count of $\mathsf T$ first reaches $f-1$ at length $f-1$ only for the
word $\mathsf T^{f-1}$ (one leaf), and at length $f$ iff $v_1\cdots v_{f-1}$ contains exactly $f-2$
letters $\mathsf T$ (one $\mathsf F$ in one of $f-1$ positions) and $v_f=\mathsf T$ (a decision, or
forced when $f=2$): $f-1$ leaves. Total: $f$ leaves, and each is a distinct node.

\emph{Row-end non-leaf nodes.} A word $w\in\{\mathsf T,\mathsf F\}^f$ with $j$ letters $\mathsf T$,
$1\le j\le f-2$, is realized by exactly one node (no prefix had $f-1$ letters $\mathsf T$), and
nothing else is. By Lemma~\ref{lem:pig2} this node is the row-start state of row $i{+}1$ with
$T_i=\{c_k:w_k=\mathsf T\}$ and $f-j\ge2$ free columns, with $\binom fj$ such words for each $j$.
By induction it has $L(f-j)$ leaves below it. Summing gives the recurrence (for $f=2$ the sum is empty
and $L(2)=2$).
\end{proof}

By the discussion of the root, $C(n)=L(n)$ for $n\ge2$, and $C(1)=1=L(1)$.

\subsection{Solving the recurrence}
We show $L(f)=f\cdot\Fub(f-1)$ for all $f\ge1$. For $f=1$: $1=1\cdot\Fub(0)$. Let $f\ge2$ and assume the claim
for smaller arguments. Using $\binom fj(f-j)=f\binom{f-1}j$,
\[
L(f)=f+\sum_{j=1}^{f-2}\binom fj (f-j)\Fub(f-j-1)=f\Bigl[1+\sum_{j=1}^{f-2}\binom{f-1}j\Fub(f-1-j)\Bigr].
\]
The Fubini recurrence for $m=f-1$ reads $\Fub(f-1)=\sum_{j=1}^{f-1}\binom{f-1}j\Fub(f-1-j)$, whose last
term ($j=f-1$) equals $\Fub(0)=1$. So the bracket is $\Fub(f-1)$ and $L(f)=f\,\Fub(f-1)$. This
proves Theorem~\ref{thm:pig}. \qed

The values are $1,2,9,52,375,3246,32781,378344,\dots$ (Table~\ref{tab:vals}). The proof was
cross-checked against a clause-level implementation of the procedure ($n\le6$), against an
independent state-based dynamic program ($n\le9$) and the recurrence of Lemma~\ref{lem:pig3} for $f\le300$
(Section~\ref{sec:repro}).

%=====================================================================
\section{Hole-major order}\label{sec:hole}

\begin{theorem}[hole-major]\label{thm:hole}
For the hole-major order, $C(1)=1$ and for every $n\ge2$
\[
C(n)=L(n):=(n+2)^{\,n-2}+(n+3)^{\,n-2}.
\]
\end{theorem}

The variables are ordered column by column: $x_{0,1},x_{1,1},\dots,x_{n,1},x_{0,2},\dots,x_{n,n}$.
For $n=1$ the root is a conflict leaf as in Section~\ref{sec:pig}. Let $n\ge2$.

\subsection{The first $n-2$ columns contain no conflict}
Say that a node is \emph{at the start of column $c$} if it is closed (propagated, conflict free),
columns $<c$ are completely assigned, and no variable of a column $\ge c$ is assigned. Let $T$ be the set of
pigeons having a true cell in columns $<c$ and $U=\{0,\dots,n\}\setminus T$.

\begin{lemma}\label{lem:h1}
For a node at the start of column $c$, $1\le c\le n-1$: every pigeon in $U$ has all its cells in columns $<c$
false, hence exactly the $n-c+1\ge2$ cells of columns $c,\dots,n$ unassigned; the pigeon clause of
every pigeon in $T$ is satisfied. The only clauses not yet satisfied are the pigeon clauses of
pigeons in $U$ and the hole clauses of columns $\ge c$.
\end{lemma}
\begin{proof}
Pigeons in $T$ satisfy their clause by definition. A pigeon in $U$ has no true cell in columns $<c$, and
those cells are assigned, hence false. A hole clause of a column $<c$
has both variables assigned and, the node being conflict free, at most one of them true, hence is satisfied.
\end{proof}

\begin{lemma}[one column, $c\le n-2$]\label{lem:h2}
Take a node at the start of column $c\le n-2$. The solver branches successively on
$x_{0,c},x_{1,c},\dots,x_{n,c}$. Then:
\begin{enumerate}
\item for each $i\in\{0,\dots,n\}$ the child ``$x_{0,c},\dots,x_{i-1,c}$ false, $x_{i,c}$ true'' is
  closed and conflict free; its assignment in column $c$ is $x_{i,c}=1$ and all other cells of column
  $c$ equal $0$, and no variable of a column $>c$ is assigned;
\item the child ``$x_{0,c},\dots,x_{n,c}$ all false'' is closed and conflict free, with column $c$
  entirely false and nothing else new assigned;
\item there are no other children and no conflict leaf occurs in column $c$.
\end{enumerate}
\end{lemma}
\begin{proof}
The only clause containing $x_{j,c}$ positively is $P_j$. Falsifying $x_{j,c}$ (by decision or by
propagation): if $j\in T$, $P_j$ is satisfied; if $j\in U$, by Lemma~\ref{lem:h1} $P_j$ has
$n-c+1\ge3$ unassigned literals (here $c\le n-2$ is used), so after falsifying this single cell of the
pigeon (the other cells of $P_j$ lie in other columns) it still has $\ge2$ unassigned literals and no
true literal: it is neither unit nor falsified. Hole clauses containing $\neg x_{j,c}$ become satisfied.
Hence setting cells of column $c$ to false triggers nothing. Setting $x_{i,c}$ true satisfies $P_i$ and
makes the hole clauses $(\neg x_{i,c}\vee\neg x_{k,c})$ unit for those $k>i$ whose cell is not yet
false (for $k<i$ the cell is already false by the earlier decisions): propagation sets $x_{k,c}=0$,
which, by what was just shown, triggers nothing further; hole clauses of other columns involve no
variable of column $c$. This gives (1) and (2); after $x_{i,c}=1$ the whole column is assigned, so the
next decision is in column $c{+}1$, which gives (3).
\end{proof}

\begin{corollary}\label{cor:h3}
Put $m=n-2$ and $A=\{\text{none},0,1,\dots,n\}$, $|A|=n+2$. The nodes at the start of column $n-1$ are in
bijection with the words $w=(w_1,\dots,w_m)\in A^m$, where $w_c$ records the outcome of column $c$
(``none'': all cells false; $i$: the first true cell is in pigeon $i$, all other cells of the column
false). All $(n+2)^m$ words occur, no conflict leaf occurs in columns $1,\dots,m$, and the node of $w$
has $T=T(w)=\{\text{pigeons occurring in }w\}$.
\end{corollary}
\begin{proof}
The root is at the start of column $1$: there are no units since $n\ge2$ (pigeon clauses have $\ge2$
literals, hole clauses are binary). Induct on $c$ using Lemma~\ref{lem:h2}: in each column all $n+2$
outcomes are available, independently of $T$, because there are no functional clauses (a pigeon with a
true cell can have another).
\end{proof}
Since $|T(w)|\le m=n-2$ we have $|U|\ge3$.

\subsection{The last two columns}
\begin{lemma}\label{lem:h4}
Let a node be at the start of column $c=n-1$ (with $d=n$ the last column) and $|U|\ge3$; let
$a=\min U$ (so the pigeons $0,\dots,a-1$ lie in $T$). The subtree below the node has exactly $a+2$
conflict leaves: the children ``$x_{0,c},\dots,x_{i-1,c}$ false, $x_{i,c}$ true'' for $i<a$ ($a$ leaves),
the child ``$x_{0,c},\dots,x_{a-1,c}$ false, $x_{a,c}$ true'' (one leaf) and the child
``$x_{0,c},\dots,x_{a,c}$ false'' (one leaf). The only non-leaf nodes are the prefixes
``$x_{0,c},\dots,x_{j,c}$ false'', $j<a$.
\end{lemma}
\begin{proof}
By Lemma~\ref{lem:h1} each $p\in U$ has exactly the two unassigned cells $x_{p,c},x_{p,d}$.
\emph{False prefixes.} For $j<a$, pigeon $j\in T$, so $P_j$ is satisfied and falsifying $x_{j,c}$
triggers nothing: no conflict, the solver goes on to $x_{j+1,c}$.
\emph{(i)} $x_{i,c}=1$, $i<a$. Propagation sets $x_{p,c}=0$ for all $p\ne i$, in particular for all
$p\in U$ ($i\in T$). Then each $P_p$, $p\in U$, becomes the unit $x_{p,d}$, so $x_{p,d}=1$ for all
$p\in U$; as $|U|\ge2$ two of them falsify a hole clause of column $d$: conflict.
\emph{(ii)} $x_{a,c}=1$. Propagation sets $x_{p,c}=0$ for $p\ne a$, in particular for the $\ge2$ pigeons of
$U\setminus\{a\}$, whose pigeon clauses become the units $x_{p,d}$, and two of them falsify a hole clause
of column $d$: conflict.
\emph{(iii)} $x_{a,c}=0$ after $x_{0,c},\dots,x_{a-1,c}$ false. Now $P_a$ is the unit $x_{a,d}$, which is
set true; the hole clauses of column $d$ force $x_{p,d}=0$ for $p\ne a$, in particular for the $\ge2$
pigeons of $U\setminus\{a\}$, whose pigeon clauses become the units $x_{p,c}=1$, and two of them
falsify a hole clause of column $c$: conflict.
These children exhaust the branching through $x_{a,c}$; (iii) ends the false path.
\end{proof}

\subsection{The exact sum}
By Corollary~\ref{cor:h3} and Lemma~\ref{lem:h4} (with $a=\mex(w)$, the least pigeon index not occurring in
$w$; it exists because $w$ has only $m<n+1$ letters),
\begin{equation}\label{eq:sum}
C(n)=\sum_{w\in A^m}\bigl(2+\mex(w)\bigr).
\end{equation}

\subsection{Evaluation of the sum}
Let $q=n+2$, $g(x)=x^m$, $(Eg)(x)=g(x+1)$ and $(\nabla g)(x)=g(x)-g(x-1)=(1-E^{-1})g(x)$; these operators
commute and $\nabla^kg(x)=\sum_{i=0}^k(-1)^i\binom ki g(x-i)$.

\emph{Step 1.} $\mex(w)=\sum_{k=1}^{n}[\text{the letters }0,\dots,k-1\text{ all occur in }w]$, so
$\sum_w\mex(w)=\sum_{k=1}^nN(k)$ with $N(k)$ the number of words in $A^m$ containing each of the $k$ fixed
letters $0,\dots,k-1$.

\emph{Step 2 (inclusion--exclusion).} $N(k)=\sum_{i=0}^k(-1)^i\binom ki(q-i)^m=(\nabla^kg)(q)$. For
$k>m$ a word of length $m$ cannot contain $k$ distinct letters, so $N(k)=0$, consistent with
$\nabla^kg=0$ for $k>m=\deg g$. Hence $\sum_w\mex(w)=\sum_{k=1}^m(\nabla^kg)(q)$.

\emph{Step 3 (identity).} For every polynomial $g$ of degree $\le m$,
\[
\sum_{k=1}^m\nabla^kg=(E-1)g.
\]
Indeed $1-\nabla=E^{-1}$ gives $\bigl(\sum_{k=1}^m\nabla^k\bigr)(1-\nabla)=\nabla-\nabla^{m+1}$. Multiplying
by $E$ and using $E\nabla=E-1$ we get $\sum_{k=1}^m\nabla^k=(E-1)-E\nabla^{m+1}$, and $\nabla^{m+1}g=0$
because each $\nabla$ lowers the degree by one and kills constants. (For $m=0$ both sides are $0$.)

\emph{Step 4.} By \eqref{eq:sum}, $C(n)=2g(q)+\bigl(g(q+1)-g(q)\bigr)=g(q)+g(q+1)=(n+2)^{n-2}+(n+3)^{n-2}$.
For $n=2$ ($m=0$, $g=1$, empty sum) this gives $C(2)=2$. This proves Theorem~\ref{thm:hole}. \qed

The values are $1,2,11,85,855,10657,159049,2771561,\dots$ (Table~\ref{tab:vals}). Verification: clause-level
procedure $n\le6$, independent state dynamic program $n\le9$, an explicit model of Lemmas~\ref{lem:h2} and
\ref{lem:h4} on subsets of pigeons ($n\le10$), direct enumeration of the word sum \eqref{eq:sum}
($n\le6$), the inclusion--exclusion form and the closed form ($n\le60$), and the identity of Step~3 on
random polynomials.

%=====================================================================
\section{Computations for other rules}\label{sec:comp}

\subsection{Exact values}
Table~\ref{tab:vals} lists the conflict-leaf counts for all four rules for $n\le14$; the anti-diagonal order
continues with
\[
C_{\mathrm{diag}}(15)=547848248627841,\qquad C_{\mathrm{diag}}(16)=15045475165081609.
\]
The anti-diagonal values for $n=1,\dots,16$ are
$1,\allowbreak 2,\allowbreak 8,\allowbreak 45,\allowbreak 331,\allowbreak 3057,\allowbreak 34007,\allowbreak 441515,\allowbreak 6541557,\allowbreak 108804265,\allowbreak 2006371441,\allowbreak 40615858425,\allowbreak 895407989033,\allowbreak 21354721047401,\allowbreak 547848248627841,\allowbreak 15045475165081609.$

\begin{table}[ht]
\centering
\scriptsize
\begin{tabular}{r r r r r}
$n$ & pigeon-major & hole-major & anti-diagonal & $\mathcal D_2$-JW \\ \hline
1 & 1 & 1 & 1 & 1 \\
2 & 2 & 2 & 2 & 4 \\
3 & 9 & 11 & 8 & 13 \\
4 & 52 & 85 & 45 & 61 \\
5 & 375 & 855 & 331 & 307 \\
6 & 3\,246 & 10\,657 & 3\,057 & 2\,110 \\
7 & 32\,781 & 159\,049 & 34\,007 & 14\,746 \\
8 & 378\,344 & 2\,771\,561 & 441\,515 & 130\,687 \\
9 & 4\,912\,515 & 55\,318\,979 & 6\,541\,557 & 1\,187\,971 \\
10 & 70\,872\,610 & 1\,245\,712\,417 & 108\,804\,265 & 12\,801\,700 \\
11 & 1\,124\,723\,193 & 31\,265\,546\,157 & 2\,006\,371\,441 & 143\,164\,903 \\
12 & 19\,471\,590\,876 & 865\,905\,045\,601 & 40\,615\,858\,425 & 1\,819\,083\,835 \\
13 & 365\,190\,378\,735 & 26\,241\,941\,903\,791 & 895\,407\,989\,033 & 24\,104\,576\,248 \\
14 & 7\,376\,016\,877\,334 & 864\,097\,213\,940\,417 & 21\,354\,721\,047\,401 & 353\,344\,753\,156 \\
\end{tabular}
\caption{Conflict leaves of DPLL with unit propagation on $\PHP(n{+}1,n)$. The first two columns are given by
Theorems~\ref{thm:pig} and~\ref{thm:hole}; the last two columns are computed (see text).}
\label{tab:vals}
\end{table}

\subsection{The anti-diagonal order}
This is the static order defined in the introduction. We know no closed form, recurrence or asymptotic law for it
and make no claim about its generating function (in particular none about D-finiteness).
The values were computed with a layered dynamic program over states $(\mathrm{pos},U,O)$, where $\mathrm{pos}$ is the
position of the next undecided variable, $U$ the set of pigeons without a true cell and $O$ the set of holes with a
true cell (program \texttt{ref\_dp.rs}). We did not prove the correctness of this reduction; instead it
was validated against (i) our clause-level implementation of the procedure for $n\le10$ and (ii) an independently
written Python port of the state dynamic program for $n\le9$. The Rust program was re-run in the preparation of this
paper for $n\le13$; the values for $n=14,15,16$ come from a single earlier run (about $29$, $137$ and $886$ seconds; $1.9\cdot10^{9}$
states at $n=16$).

\subsection{The rule $\mathcal D_2$-JW}\label{sec:res2}
\begin{definition}[$\mathcal D_2$-JW]
Given a node with partial assignment $\alpha$, run unit propagation; on a conflict the node is a leaf. Otherwise,
for every unassigned variable $x$ define its Jeroslow--Wang score \cite{JW90}
\[
J(x)=\sum_{C\ni x\text{ or }\neg x,\ C\text{ not satisfied by }\alpha}2^{-u(C)},
\]
where $u(C)$ is the number of unassigned literals of $C$ (both polarities of $x$ are summed). Let $x^{(1)},x^{(2)}$ be
the two unassigned variables with the largest scores, ties broken in favour of the variable that is earlier in
row-major order ($x_{0,1},x_{0,2},\dots$) (by Proposition~\ref{prop:mod3} below, at least six variables are unassigned at every non-conflict node, so two variables can always be chosen).
The node has four children, one for each of the four assignments of $(x^{(1)},x^{(2)})$, each followed by unit
propagation. $r(n)$ is the number of conflict leaves.
\end{definition}
(For $\PHP$ the clauses containing $x_{i,j}$ that are not satisfied are $P_i$, if $i$ has no true cell, and the
hole clauses $\{x_{i,j},x_{k,j}\}$ with $x_{k,j}$ unassigned, so $J(x_{i,j})=2^{-u(P_i)}[P_i\text{ open}]+(u_j-1)/4$ with $u_j$ the number of unassigned
cells in column $j$.) The values are $r(1),\dots,r(14)=1,\allowbreak 4,\allowbreak 13,\allowbreak 61,\allowbreak 307,\allowbreak 2110,\allowbreak 14746,\allowbreak 130687,\allowbreak 1187971,\allowbreak 12801700,\allowbreak 143164903,\allowbreak 1819083835,\allowbreak 24104576248,\allowbreak 353344753156$.
They were computed by a memoised search on residual states canonical under order-preserving renaming of holes and
rows (program \texttt{ref\_resk.rs}); the number of memoised states is $1127$ for $n=8$ and $10\,884\,108$ for $n=14$. We verified the
state-based search against a direct clause-level implementation of the definition for $n\le7$ ($k=2$; $n\le7$ also for $k=3,4$ in the companion logs) and against an independent
Python port for $n\le12$. The Rust program was re-run for $n\le12$ here; $r(13)$ (about one minute) and $r(14)$ (about one hour)
are taken from earlier runs; in a later independent audit, $r(1),\dots,r(15)$ were also reproduced by separately written programs (a state-compressed Rust program for $n\le15$), whereas $r(16)$ has not been reproduced. An earlier audit found the same counts for $n\le 8$ when ties are broken towards the \emph{last} variable instead
(not re-run here). We know no closed form for $r(n)$, and no proof of any growth law; the congruence $r(n)\equiv1\pmod3$ is proved in Proposition~\ref{prop:mod3}.

\begin{proposition}\label{prop:mod3}
Let $n\ge2$. Every partial assignment of $\PHP(n{+}1,n)$ that falsifies no clause and has no unit clause leaves at least six variables unassigned. Consequently every internal node of the $\mathcal D_2$-JW tree has exactly four children, $r(n)=3I(n)+1$, and $r(n)\equiv1\pmod 3$ for all $n\ge1$. The same holds, with $2^k$ children and modulus $2^k-1$, for the $k$-variable rule for $2\le k\le6$, and for any rule that chooses $k$ unassigned variables.
\end{proposition}
\begin{proof}
Let $T$ be the pigeons with a true cell and $U$ the others. Distinct pigeons of $T$ have true cells in distinct holes (a hole clause would be falsified), so the set $O$ of holes with a true cell has $|O|\ge|T|$; put $e=n-|O|$. Then $|U|=n+1-|T|\ge e+1\ge1$. For $p\in U$ and $h\in O$ the hole clause with the true cell of hole $h$ would be unit or falsified unless $x_{p,h}$ is false; so all unassigned cells of row $p$ lie in the $e$ holes outside $O$. Since $p\in U$, $P_p$ has no true literal; it is neither falsified nor unit, so it has at least two unassigned literals; hence $e\ge2$, $|U|\ge3$, and the rows of $U$ contribute at least $3\cdot2=6$ unassigned variables. At a non-conflict node (closed assignment, Lemma~\ref{lem:conf}) the lemma applies, so $\min(k,m)=k$ for $k\le6$ and each internal node has $2^k$ children; for $n=1$ the root is a conflict leaf. A tree in which every internal node has $2^k$ children satisfies $L=(2^k-1)I+1$ by edge counting.
\end{proof}

\subsection{Reproducibility}\label{sec:repro}
The files \texttt{ref\_dp.rs} (static-order dynamic program; usage
\texttt{ref\_dp <pig|hole|diag> nmin nmax threads}), \texttt{ref\_resk.rs} (Res($k$)-JW; usage
\texttt{ref\_resk k nmin nmax}) and \texttt{ms\_check.py} accompany the paper and are deposited at
\url{https://github.com/arunkatherashala/php-dpll-counts} (directory \texttt{code/}). \texttt{ms\_check.py} recomputes, in exact integer or rational arithmetic, the
clause-level counts, the state dynamic programs, all closed forms for $n\le60$ (and the recurrence of
Lemma~\ref{lem:pig3} for $f\le300$), the identities used in the proofs, the counterexamples of
Section~\ref{sec:errata}, and the certificates of Theorem~\ref{thm:C}; it runs the Rust executables when present
and prints \texttt{ALL OK} at the end. Items it does not recompute (the stored runs mentioned above and the
$n=4$ minimum tree-resolution value, available with \texttt{--slow}) are listed at the end of its output.

\paragraph{Machine-checked parts.} The directory \texttt{lean/} of the same repository contains Lean~4 files
(core library only, no Mathlib, no \texttt{sorry}; the only axioms used are the standard ones \texttt{propext},
\texttt{Quot.sound} and \texttt{Classical.choice}). They verify, for \emph{all} $f$ and $n$, the purely combinatorial
halves of Theorems~\ref{thm:pig} and~\ref{thm:hole}: that the recurrence of Lemma~\ref{lem:pig3} has the solution
$L(f)=f\cdot\Fub(f-1)$ (\texttt{Pigeon.lean}), and that the word sum \eqref{eq:sum} equals
$(n+2)^{n-2}+(n+3)^{n-2}$ for $n\ge2$ (\texttt{Hole.lean}, \texttt{Hole2.lean}). The reductions from the DPLL
procedure with unit propagation to the recurrence and to the word sum (Lemmas~\ref{lem:pig1}, \ref{lem:pig2},
\ref{lem:h1}, \ref{lem:h2} and~\ref{lem:h4}, Corollary~\ref{cor:h3}) are \emph{not} formalized; they rest on the proofs given here.

\subsection{Earlier claims that are false or not established}\label{sec:errata}
Earlier drafts of this work (a series of manuscripts on the same numerical data) contained statements that we
withdraw or correct here. We list them, so that nothing in this paper is read as supporting them.
\begin{enumerate}
\item \emph{``These are the first exact formulas for tree-resolution proof size.''} Misdescribed. The
  quantities are DPLL-with-UP conflict counts for a fixed branching heuristic. The true minimum number of
  leaves of tree-like resolution refutations of $\PHP(n{+}1,n)$ is $3,11,43,189$ for $n=1,\dots,4$
  (recomputed here by exhaustive search over decision trees for $n\le3$; the value $189$ for $n=4$ was obtained in an
  earlier audit and by a run of the same search, available with \texttt{ms\_check.py --slow}), compared with
  pigeon-major counts $1,2,9,52$.
\item \emph{``The dynamic Jeroslow--Wang heuristic gives exactly $n!$, and $n!$ is the minimum over all
  orderings.''} The minimality statement is false for adaptive variable choice: minimizing over \emph{all} adaptive
  choices of the branching variable at every node, the minimal number of conflict leaves of DPLL with
  unit propagation is $1,2,6,23$ for $n=1,\dots,4$ (exhaustive search over canonical states), and $23<4!=24$.
  Whether some \emph{static} order beats $n!$ at $n=4$ is not decided here. The claimed value $n!$ for the dynamic rule is not
  established in this paper.
\item \emph{Asymptotics of the pigeon-major counts.} It was claimed that $A052882(n)\sim n!/(2\ln2)^n$, i.e.\ that
  $C(n)/n!\to0$ geometrically. This is wrong: $C(n)/n!\sim\tfrac12(\ln2)^{-n}$ grows (Proposition~\ref{prop:pigasym}). A
  claim that the hole-major count divided by $n!$ grows like $2.47^n$, with the EGF coefficient ratio converging to $e-\tfrac12\approx2.22$, is also wrong; the
  ratio of consecutive EGF coefficients tends to $e$ (Proposition~\ref{prop:holeasym}).
\item \emph{``The diagonal sequence is not D-finite''} (with the associated Pad\'e/ODE ``asymptotic law'',
  $\rho\approx0.544$, etc.). Not established. The argument consisted of rank/least-squares tests on a finite prefix of the
  sequence; such a test can at most show that a prefix is not annihilated by an ODE of a given small order and degree (and
  every finite prefix is annihilated by some larger ODE). Furthermore the text defined the sequence by a diagonal
  order with Jeroslow--Wang scoring, whereas the tabulated numbers are those of the static anti-diagonal order
  defined here. We have no result on the type of singularity of the anti-diagonal generating function.
\item \emph{A ``crossover theorem'': $\mathcal D_2$-JW has the smallest tree among four orderings for all $n\ge10$.} Not a
  theorem, and not supported as stated. With the exact values of Table~\ref{tab:vals}, $\mathcal D_2$-JW is the smallest of the four counts for every
  $5\le n\le14$, so there is no crossover at $n=10$ in this data (the earlier table contained inconsistent pigeon-major and hole-major columns).
  The four numbers also count leaves of trees with different branching (binary versus 4-way) and
  different rules, so a comparison is not a statement about problem difficulty.
\item \emph{A quasi-polynomial closed form} $a(n)=(n^2+13n+c(n\bmod6))/6$, $c=(6,10,6,12,4,12)$, for $a(n)=\lfloor r(n{+}1)/n!\rfloor$,
  and the conclusions drawn from it (a triple pole, and so on). False. $a(n)=1,4,6,10,12,17,20,25,29,35,39$ for $n\le10$ agrees with the formula, but
  $a(11)=45$ (since $r(12)/11!=45.5718\dots$) whereas the formula gives $46$; also $a(12)=50$ versus $51$ and $a(13)=56$ versus $58$. Eleven
  data points do not support a formula with six free periodic parameters.
\item \emph{Res(2) singularity claims} (a singularity at $\rho\approx0.894$, exponent $3/2$): numerical conjectures that we have neither
  re-derived nor used.
\item \emph{Logarithmic singularity ``proved computationally'' from 12 terms.} Finite data cannot establish a statement about the
  radius of convergence of an infinite series (Remark~\ref{rem:tail}); the correct status is Theorem~\ref{thm:C}, which is conditional.
  In that line of work the printed recurrence for the coefficients of $\log R$ must read $nf_n=nb_n-\sum_{k<n}kf_kb_{n-k}$, and the claim
  $F_n\equiv0\pmod 6$ for $n\ge4$ is false ($F_6=-31059$ is odd).
\end{enumerate}

%=====================================================================
\section{Generating functions and asymptotics}\label{sec:asym}

\subsection{Pigeon-major}
By Theorem~\ref{thm:pig}, $\sum_{n\ge1}C(n)\,x^n/n!=\sum_{n\ge1}\Fub(n-1)x^n/(n-1)!=x/(2-e^x)$.

\begin{proposition}\label{prop:pigasym}
$\displaystyle C(n)=\frac{n!}{2(\ln2)^n}\bigl(1+O(\kappa^n)\bigr)$ with $\kappa=\ln2/|\ln2+2\pi i|=0.1097\ldots$. More precisely,
for $n\ge2$ one has the pole expansion $C(n)/n!=\tfrac12\sum_{k\in\mathbb Z}(\ln2+2\pi ik)^{-n}$. Equivalently
$\Fub(m)\sim m!/(2(\ln2)^{m+1})$ and $C(n)=n\Fub(n-1)$.
\end{proposition}
\begin{proof}
$x/(2-e^x)$ is meromorphic with simple poles at $x_k=\ln2+2\pi ik$, $k\in\mathbb Z$; near $x_k$,
$2-e^x=-2(x-x_k)+O((x-x_k)^2)$, so the principal part is $\tfrac{x_k}{2(x_k-x)}$, whose $n$-th coefficient is $\tfrac12x_k^{-n}$. The
dominant singularity is $x_0=\ln2$; the next ones have modulus $|x_{\pm1}|=6.32\ldots$. The asymptotic statement is the standard
theorem for meromorphic functions with a unique dominant simple pole \cite[Thm.~IV.10]{FS09}. The full
pole expansion follows by the same contour argument and was checked numerically (Section~\ref{sec:repro}).
\end{proof}
Numerically $C(n)\,2(\ln2)^n/n!=0.99999999977$ at $n=10$ and $1-O(10^{-15})$ for $n\ge20$. In particular
$C(n)/n!\to\infty$ (like $\tfrac12\cdot1.4427^n$). Since the generating function has infinitely many poles it is not D-finite (a D-finite
function has finitely many singularities).

\subsection{Hole-major}
Let $T(x)=\sum_{n\ge1}n^{n-1}x^n/n!$ be the tree function, $T=xe^{T}$, with radius of convergence $1/e$ and the expansion
$T(x)=1-\sqrt2\,s^{1/2}+\tfrac23s-\cdots$, $s=1-ex$ at $x=1/e$.
\begin{proposition}\label{prop:holeasym}
For $a\ge1$,
\[
\sum_{n\ge0}(n+a)^{n-2}\frac{x^n}{n!}=e^{aT}\Bigl(\frac1{a^2}-\frac{T}{a(a+1)}\Bigr).
\]
Consequently $H(x):=\sum_{n\ge1}C(n)x^n/n!=e^{2T}\bigl(\tfrac14-\tfrac T6\bigr)+e^{3T}\bigl(\tfrac19-\tfrac T{12}\bigr)-\tfrac{13}{36}+\tfrac5{12}x$. $H$ has a unique
dominant singularity at $x=1/e$ of type $(1-ex)^{3/2}$:
\[
H(x)=\text{(analytic)}+\tfrac{2\sqrt2}{3}\,(e^2+e^3)(1-ex)^{3/2}+O\bigl((1-ex)^2\bigr),
\]
and
\[
\frac{C(n)}{n!}=\frac{e^2+e^3}{\sqrt{2\pi}}\,\frac{e^{n}}{n^{5/2}}\bigl(1+O(1/n)\bigr),\qquad
C(n)=n^{\,n-2}\Bigl(e^2+e^3-\frac{6e^2+\frac{21}2e^3}{n}+O(n^{-2})\Bigr).
\]
In particular $C(n)\sim(e^2+e^3)\,n^{n-2}=27.4745\ldots\,n^{n-2}$, and the ratio of consecutive coefficients $\frac{C(n+1)/(n+1)!}{C(n)/n!}$ tends to $e$.
\end{proposition}
\begin{proof}
From $T=xe^T$, Lagrange inversion gives $e^{aT}=\sum_{n\ge0}a(n+a)^{n-1}x^n/n!$, i.e.\ $T(t)^a/t=\sum_{n\ge0}a(n+a)^{n-1}t^{n+a-1}/n!$ (using $T^a=t^ae^{aT}$).
Integrating termwise and using $\int_0^xt^{n+a-1}dt=x^{n+a}/(n+a)$ we get, for $G_a(x):=\sum_{n\ge0}a(n+a)^{n-2}x^{n+a}/n!$,
$G_a(x)=\int_0^xT(t)^a\,dt/t$. Substituting $u=T(t)$, $t=ue^{-u}$, $dt/t=(1-u)\,du/u$ (valid for formal power series), we get
$G_a=\int_0^{T}u^{a-1}(1-u)\,du=T^a/a-T^{a+1}/(a+1)$. Dividing by $ax^a$ and using $T^ax^{-a}=e^{aT}$ gives the first formula. The formula for $H$ follows from $a=2,3$, because $(n+a)^{n-2}$ for $n=0,1$ gives the constant terms $\tfrac14+\tfrac19=\tfrac{13}{36}$ and
$\tfrac13+\tfrac14=\tfrac7{12}$ of the sum, and $C(1)=1$: $1-\tfrac7{12}=\tfrac5{12}$.

Singular behaviour. Put $B_a(T)=e^{aT}(a^{-2}-T/(a(a+1)))$. Then $B_a'(T)=e^{aT}b(T)$ with $b(T)=\tfrac1a-\tfrac T{a+1}-\tfrac1{a(a+1)}$, and $b(1)=0$; thus $B_a'(1)=0$ (the
square-root term cancels), $B_a''(1)=-e^a/(a+1)$, $B_a'''(1)=-2ae^a/(a+1)$. Inserting $T-1=-\sqrt2s^{1/2}+\tfrac23s+O(s^{3/2})$ gives
$(T-1)^2=2s-\tfrac{4\sqrt2}3s^{3/2}+O(s^2)$ and $(T-1)^3=-2\sqrt2s^{3/2}+O(s^2)$, so the coefficient of $s^{3/2}$ in $B_a(T(x))$ is
$-\tfrac{\sqrt2}{3}\bigl(2B_a''(1)+B_a'''(1)\bigr)=\tfrac{2\sqrt2}3e^a$. Since $T$ is analytic in a $\Delta$-domain at $1/e$ (the only singularity of $T$ on its circle of convergence) the transfer theorem
\cite[Thm.~VI.3 and Cor.~VI.1]{FS09} with the coefficients of the standard scale $[x^n](1-ex)^{3/2}\sim e^nn^{-5/2}/\Gamma(-\tfrac32)$, $\Gamma(-\tfrac32)=\tfrac{4\sqrt\pi}3$ (Thm.~VI.1), yields
$C(n)/n!\sim\tfrac{2\sqrt2}3(e^2+e^3)\tfrac3{4\sqrt\pi}\,e^nn^{-5/2}=\tfrac{e^2+e^3}{\sqrt{2\pi}}e^nn^{-5/2}$.
Finally $(1+a/n)^{n-2}=\exp\bigl((n-2)(a/n-a^2/2n^2+O(n^{-3}))\bigr)=e^a\bigl(1-(2a+a^2/2)/n+O(n^{-2})\bigr)$, giving the stated expansion of $C(n)/n^{n-2}$ with $a=2,3$.
\end{proof}
\medskip
\noindent\textbf{Numerical confirmation} (Section~\ref{sec:repro}): $C(n)/n^{n-2}=27.4695$ at $n=50\,000$ against $e^2+e^3=27.4746$; $n\bigl(C(n)/n^{n-2}-e^2-e^3\bigr)=-255.20$ at $n=50\,000$ against
$-(6e^2+10.5e^3)=-255.23$; the EGF identity of the proposition was verified as an identity of rational power series through degree $29$; the asymptotics
of $C(n)/n!$ were confirmed at $n=200,2000,20000$ to within $40/n$ in the logarithm.

\subsection{Other rules}
We have no asymptotic statement for the anti-diagonal order or for $\mathcal D_2$-JW (Section~\ref{sec:comp}).

%=====================================================================
\section{A conditional result for the $\mathcal D_2$-JW generating function}\label{sec:cond}

Let $r(n)$ be the $\mathcal D_2$-JW counts, $b_n=r(n+1)/n!$ ($n\ge0$, so $b_0=r(1)=1$), $R(x)=\sum_{n\ge0}b_nx^n$ (this is $\frac{d}{dx}\sum_{n\ge1}r(n)x^n/n!$), and let
$C(x)=\log R(x)=\sum_{n\ge1}f_nx^n$ be the formal logarithm, $F_n=n!f_n$. The $f_n$ obey $nf_n=nb_n-\sum_{k=1}^{n-1}kf_kb_{n-k}$ and, for $n\le13$, are determined by $b_0,\dots,b_n$.
(We write $C(x)$ here for the generating function of the logarithm; it has nothing to do with the leaf counts $C(n)$ of the other sections, which we do not use in this section.)
Exactly, $F_1,\dots,F_{13}=4,\allowbreak -3,\allowbreak 33,\allowbreak -216,\allowbreak 2496,\allowbreak -31059,\allowbreak 479934,\allowbreak -8454426,\allowbreak 171387720,\allowbreak -3898302102,\allowbreak 98654494668,\allowbreak -2745253754127,\allowbreak 83351925966882$.

\paragraph{The 14-term zero.} The polynomial $R_{14}(x)=\sum_{n=0}^{13}b_nx^n$ has exactly one zero in $|x|<9/20$, a real simple zero
\[
x_0^{(14)}=-0.3952047138\ldots,\qquad R_{14}'(x_0^{(14)})=1.8956,
\]
(the 12-term polynomial had $-0.394992$). The numbers $nf_nx_0^n$, $x_0=x_0^{(14)}$, for $n=7,\dots,13$ are
$-1.0037,-0.9982,-0.9997,-0.9985,-0.9986,-0.9984,-0.9983$, i.e.\ close to $-1$, the value expected if $C$ has a logarithmic singularity at $x_0$.
This is numerical evidence only: it uses nothing but $b_0,\dots,b_{13}$.

\begin{remark}\label{rem:tail}
The data cannot decide the question. Changing only the tail: $b_{14}=8899.6$ and $b_n=0$ ($n>14$) makes $-0.42$ an exact zero of $R$ (so the
dominant zero moves to $-0.42$), and $b_n=n!$ for $n\ge14$ gives a series with radius of convergence $0$. Any statement about the radius of convergence
of $\log R$ is a statement about the infinite tail of the sequence $r(n)$ and cannot follow from $r(1),\dots,r(14)$.
\end{remark}

\begin{hypothesis}[H --- \textbf{unproved}]\label{hyp:H}
For all $n\ge14$: $r(n+1)\le n^2\,n!$, i.e.\ $b_n\le n^2$.
\end{hypothesis}
Hypothesis~\ref{hyp:H} is an assumption about an infinite tail. We could not prove it: the generic bound ``leaves $\le 4^{\text{depth}}$'' gives only $r(n)\le 4^{\lceil n(n+1)/2\rceil}$,
which is not even exponentially bounded in $b_n$, and we did not find a structural (potential-function) argument for the Jeroslow--Wang dynamics. The data are consistent with it: $b_n/n^2=0.460,0.436,0.395,0.377,0.349,0.336$ for
$n=8,\dots,13$, decreasing. It is not proved, and passing 6 data points proves nothing about the tail.

\begin{theorem}[conditional on Hypothesis~\ref{hyp:H}]\label{thm:C}
Assume Hypothesis~\ref{hyp:H}. Then:
\begin{enumerate}
\item $R$ is analytic in $|x|<1$.
\item $R$ has exactly one zero $x_0$ in $|x|<9/20$ (counted with multiplicity); it is simple, real, and $-0.395638<x_0<-0.394785$; in particular $x_0\ne-2/5$. Moreover $R<0$ on $[-0.6,-0.4]$ and $R>0$ on $[0,1)$,
  so $x_0$ is the only zero of $R$ on $(-0.6,1)$. (We do \emph{not} claim uniqueness of a negative zero in $(-1,0)$.)
\item The power series of $C=\log R$ has radius of convergence exactly $|x_0|$.
\item For every $\rho<9/20$, $f_n=-x_0^{-n}/n+O(\rho^{-n})$. Hence $nf_nx_0^n\to-1$ (with error $O(n(|x_0|/\rho)^n)$ for any $\rho\in(|x_0|,9/20)$), $f_n/f_{n+1}\to x_0$, and $\operatorname{sign}f_n=(-1)^{n+1}$ for large $n$.
\end{enumerate}
\end{theorem}
\begin{proof}
Write $R_{14}=\sum_{n<14}b_nx^n$ (exact rationals). Under H, for $|x|=\rho<1$,
\[
|R(x)-R_{14}(x)|\le T(\rho):=\sum_{n\ge14}n^2\rho^n=\rho^{14}\,\frac{196-363\rho+169\rho^2}{(1-\rho)^3}.
\]
(1) follows from $0<b_n\le n^2$ for $n\ge14$.

\emph{Step A (sign change).} Exact rational arithmetic gives $R_{14}(-0.39)-T(0.39)>0$ and $R_{14}(-0.40)+T(0.40)<0$; hence $R(-0.39)>0>R(-0.40)$ and the continuous real function $R$ has a zero in $(-0.40,-0.39)$. Bisection on
the same inequalities localizes it in $(-0.395638,-0.394785)$.

\emph{Step B (zero count).} We show $|R_{14}|>T(9/20)=5.61\cdot10^{-3}$ on $|x|=9/20$. The circle is sampled at $3001$ exact rational points $z=\rho(1-t^2+2it)/(1+t^2)$ ($t$ rational, two half-circles, so $|z|=\rho$ exactly);
$L=\sum nb_n\rho^{n-1}=30.1\ge\max_{|z|=\rho}|R_{14}'|$; the arc length between consecutive samples is at most $ds=4\rho/1500$. Exact arithmetic verifies for every arc that
$\min(|w_j|,|w_{j+1}|)-L\,ds/2>T(\rho)$, where $w_j=R_{14}(z_j)$ (the minimum of $|R_{14}|$ over the samples is $0.106$). Since each point of an arc is within $ds/2$ of a sample and $R_{14}$ is $L$-Lipschitz, $|R_{14}|>T>|R-R_{14}|$ on the whole circle.
By Rouch\'e, $R$ and $R_{14}$ have the same number of zeros in $|x|<9/20$. The argument change of $R_{14}$ along each arc is the principal value $\operatorname{Arg}(w_{j+1}/w_j)$ (verified exactly: $L\,ds<|w_j|$), and these sum to $2\pi\cdot1.000000$, so $R_{14}$ has exactly one zero in the
disc. Hence so has $R$; by Step A it is real and it is simple. (The floating-point \texttt{atan2} is used only to add the angle increments, each of modulus $<\pi/2$; the total is within $10^{-6}$ of an integer.)

\emph{Step C (negative axis).} On a grid of step $10^{-4}$ in $[0.4,0.6]$ exact arithmetic verifies $R_{14}(-x_i)+L_i\,\text{step}+T(x_{i+1})<0$, where $L_i=\sum nb_nx_{i+1}^{n-1}$ bounds $|R_{14}'|$; thus $R<0$ on $[-0.6,-0.4]$. $R>0$ on $[0,1)$ because $b_n>0$. With Step B this proves (2).

\emph{Step D (radius of $\log R$).} $D_0=\{|x|<|x_0|\}$ is simply connected, $R$ has no zero in $D_0$ and $R(0)=1$, so there is an analytic $C$ on $D_0$ with $e^{C}=R$, $C(0)=0$, whose Taylor series is the formal logarithm; the radius is $\ge|x_0|$.
If it were $>|x_0|$, $C$ would extend analytically to a disc $D_1\ni x_0$ (take $D_1=\{|x|<\rho'\}$, $|x_0|<\rho'<9/20$), and $e^{C}=R$ on $D_1$ by the identity theorem, so $R(x_0)=e^{C(x_0)}\ne0$, a contradiction.

\emph{Step E (coefficients).} $h(x):=R(x)/(1-x/x_0)$ is analytic in $|x|<1$ (the simple zero cancels) and has no zero in $|x|<9/20$ by Step B. So $\log h$ ($h(0)=1$) is analytic in this disc and by Cauchy estimates
$[x^n]\log h=O(\rho^{-n})$ for every $\rho<9/20$. On $D_0$, $C=\log h+\log(1-x/x_0)$ (both sides are logarithms of $R$ vanishing at $0$) and $[x^n]\log(1-x/x_0)=-x_0^{-n}/n$. This gives (4) and, as $|x_0|<0.3957<9/20$, the consequences listed. \qedhere
\end{proof}

The certificates of Steps A--C are verified in exact rational arithmetic by \texttt{ms\_check.py}; the proof otherwise uses only standard complex analysis. The effective constants in (4) were not computed.
The theorem does \emph{not} establish that $R$, or $C$, has any property unconditionally, and it says nothing about whether Hypothesis~\ref{hyp:H} holds.

%=====================================================================
\section{Related work}\label{sec:related}

Resolution proof complexity of the pigeonhole principle starts with Haken's exponential lower bound for general resolution \cite{Haken85}; the notion of a proof system and
its relation to resolution is due to Cook and Reckhow \cite{CR79}. Buss and Tur\'an extended Haken's bound to generalized pigeonhole principles \cite{BT88}; Buss and Pitassi give upper bounds for the weak pigeonhole principle in resolution and
lower bounds for tree-like resolution, and a normal form for resolution proofs of pigeonhole principles \cite{BP98}.
The tree-like resolution size of $\PHP$ is $2^{\Theta(n\log n)}$-type: the lower bound $2^{\Omega(n\log n)}$ is due to Iwama and Miyazaki \cite{IM99} and independently to Dantchev and Riis \cite{DR01}, with a short proof by Beyersdorff, Galesi and Lauria via asymmetric Prover--Delayer games \cite{BGL10}. Urquhart \cite{Urq87} proved exponential lower bounds for Tseitin formulas on expanders. Resolution (DAG-like) refutations of size $n^32^n$ are given in \cite{BP98}.
The search procedure studied here is the DPLL procedure with unit propagation \cite{DP60,DLL62}, with the Jeroslow--Wang scoring rule \cite{JW90}. The sequence $n\Fub(n-1)$ is listed as A052882 in the OEIS \cite{OEIS}
(A052882 and A000670 were re-read in the live OEIS on 2026-10-03: $a(n)=n\cdot\mathrm{A000670}(n-1)$ and $a(n)\sim n!/(2(\ln2)^n)$); in an earlier search of the OEIS (run by us, not re-run), the hole-major sequence $1,2,11,85,855,\dots$, the anti-diagonal sequence
and the $\mathcal D_2$-JW sequence did not appear (this search was not repeated). Generating-function and singularity-analysis facts are from Flajolet and Sedgewick \cite{FS09}.
We are not aware of a published exact enumeration of the conflict leaves of DPLL with unit propagation on the pigeonhole formula for specific static orders, but we did not perform a systematic search, so ``to our knowledge'' applies.

%=====================================================================
\section{Limitations}\label{sec:limits}
\begin{itemize}
\item The results are about one specific, weak procedure (unit propagation only, binary branching on the first unassigned variable of a fixed order) on one specific formulation of $\PHP$ (no functional, no onto clauses).
  They are sensitive to the order: by Theorems~\ref{thm:pig} and~\ref{thm:hole} and Section~\ref{sec:asym}, the hole-major count exceeds the pigeon-major count by a factor of order $(e\ln2)^n n^{-5/2}$, i.e.\ exponentially in $n$. They say nothing about CDCL solvers, about other encodings, or about proof size.
\item They are not proof-complexity bounds: see the discussion in Section~\ref{sec:intro}, where the minimal tree-like resolution sizes $3,11,43,189$ differ from the counts.
\item Theorems~\ref{thm:pig} and~\ref{thm:hole} are proved here by hand. The proofs have been checked against traces and counts for small $n$ by computer, but have not been formalized in a proof assistant or independently refereed.
\item The anti-diagonal values for $n>12$ and the $\mathcal D_2$-JW values for $n>12$ rest on one Rust implementation each (cross-validated against independent implementations up to $n=12$); $C_{\mathrm{diag}}(16)$ was produced once and not re-run (and $r(16)$, if used, is likewise unconfirmed); $r(14)$ and $C_{\mathrm{diag}}(14),C_{\mathrm{diag}}(15)$ were later reproduced by independently written programs. The
  $\mathcal D_2$-JW values depend on the precise definition (scoring, tie-breaking, treatment of satisfied clauses) given in Section~\ref{sec:res2}.
\item No closed form, recurrence, or asymptotic law is known to us for the anti-diagonal or $\mathcal D_2$-JW counts; nothing is proved about D-finiteness or any other class membership for them. The congruence $r(n)\equiv1\pmod3$ is proved (Proposition~\ref{prop:mod3}).
\item Theorem~\ref{thm:C} is conditional on Hypothesis~\ref{hyp:H}, which is unproved, and its core certificates are computer-assisted (exact rational arithmetic with a floating-point summation of angle increments with exact margins).
\item The literature check is limited to the sources listed in the bibliography; the OEIS search for the hole-major, anti-diagonal and $\mathcal D_2$-JW sequences was not repeated.
\end{itemize}

%=====================================================================
\begin{thebibliography}{99}
\bibitem{CR79} S.~A. Cook and R.~A. Reckhow, The relative efficiency of propositional proof systems, \emph{J. Symbolic Logic} \textbf{44}(1) (1979), 36--50, doi:10.2307/2273702.
\bibitem{DP60} M.~Davis and H.~Putnam, A computing procedure for quantification theory, \emph{J. ACM} \textbf{7}(3) (1960), 201--215, doi:10.1145/321033.321034.
\bibitem{DLL62} M.~Davis, G.~Logemann and D.~Loveland, A machine program for theorem-proving, \emph{Comm. ACM} \textbf{5}(7) (1962), 394--397, doi:10.1145/368273.368557.
\bibitem{FS09} P.~Flajolet and R.~Sedgewick, \emph{Analytic Combinatorics}, Cambridge University Press, 2009, doi:10.1017/CBO9780511801655.
\bibitem{Haken85} A.~Haken, The intractability of resolution, \emph{Theoret. Comput. Sci.} \textbf{39} (1985), 297--308, doi:10.1016/0304-3975(85)90144-6.
\bibitem{JW90} R.~G. Jeroslow and J.~Wang, Solving propositional satisfiability problems, \emph{Ann. Math. Artif. Intell.} \textbf{1}(1--4) (1990), 167--187, doi:10.1007/BF01531077.
\bibitem{BP98} S.~R. Buss and T.~Pitassi, Resolution and the weak pigeonhole principle, in \emph{Computer Science Logic (CSL'97)}, LNCS \textbf{1414}, Springer, 1998, 149--156, doi:10.1007/BFb0028012.
\bibitem{BT88} S.~R. Buss and G.~Tur\'an, Resolution proofs of generalized pigeonhole principles, \emph{Theoret. Comput. Sci.} \textbf{62} (1988), 311--317, doi:10.1016/0304-3975(88)90072-2.
\bibitem{IM99} K.~Iwama and S.~Miyazaki, Tree-like resolution is superpolynomially slower than DAG-like resolution for the pigeonhole principle, in \emph{ISAAC '99}, LNCS \textbf{1741}, Springer, 1999, 133--142, doi:10.1007/3-540-46632-0\_14.
\bibitem{DR01} S.~S. Dantchev and S.~Riis, Tree resolution proofs of the weak pigeon-hole principle, in \emph{CCC 2001}, IEEE, 2001, 69--75, doi:10.1109/CCC.2001.933873.
\bibitem{BGL10} O.~Beyersdorff, N.~Galesi and M.~Lauria, A lower bound for the pigeonhole principle in tree-like resolution by asymmetric Prover--Delayer games, \emph{Inform. Process. Lett.} \textbf{110}(23) (2010), 1074--1077, doi:10.1016/j.ipl.2010.09.007.
\bibitem{Urq87} A.~Urquhart, Hard examples for resolution, \emph{J. ACM} \textbf{34}(1) (1987), 209--219, doi:10.1145/7531.8928.
\bibitem{OEIS} The On-Line Encyclopedia of Integer Sequences, sequences A052882 and A000670, \texttt{https://oeis.org}.
\end{thebibliography}

\end{document}
